%% General definitions
\documentclass{article} %% Determines the general format.
\usepackage{a4wide} %% paper size: A4.
\usepackage[utf8]{inputenc} %% This file is written in UTF-8.
%% Some editors on Windows cannot save files in UTF-8.
%% If there is a problem with special characters not showing up
%% correctly, try switching "utf8" to "latin1" (ISO 8859-1).
\usepackage[T1]{fontenc} %% Format of the resulting PDF file.
\usepackage{fancyhdr} %% Package to create a header on each page.
\usepackage{lastpage} %% Used for "Page X of Y" in the header.
                      %% For this to work, you have to call pdflatex twice.
\usepackage{enumerate} %% Used to change the style of enumerations (see below).

\usepackage{amssymb} %% Definitions for math symbols.
\usepackage{amsmath} %% Definitions for math symbols.

\usepackage{tikz}  %% Pagacke to create graphics (graphs, automata, etc.)
\usetikzlibrary{automata} %% Tikz library to draw automata
\usetikzlibrary{arrows}   %% Tikz library for nicer arrow heads


%% Left side of header
\lhead{\course\\\semester\\Exercise \homeworkNumber}
%% Right side of header
\rhead{\authorname\\Page \thepage\ of \pageref{LastPage}}
%% Height of header
\usepackage[headheight=36pt]{geometry}
%% Page style that uses the header
\pagestyle{fancy}

\newcommand{\authorname}{Luis Tritschler\\Benjamin Regez}
\newcommand{\semester}{Spring Semester 2026}
\newcommand{\course}{Foundations of Artificial Intelligence}
\newcommand{\homeworkNumber}{1}


\begin{document}

\section*{Statement on scientific integrity}
Google Translate was used to assist in translating individual words from
German to English when writing down the exercise solutions.\\
Apart from that, we solved the entire exercise sheet on our own. 

\section*{Exercise 1.1}

\begin{enumerate}[(a)]
    \item Pepper categorizes human emotions based on knowledge about facial
    expressions and voice tones. This requires knowledge about how humans act,
    therefore, Pepper belongs to the category of \texttt{acting humanly}.
    \item AlphaGo decides the best move based on knowledge from games in the past,
    both played by humans and computers. We would therefore argue that AlphaGo most
    likely can be seen as an \texttt{acting rationally} agent, especially because choosing
    moves with the greatest probatility of winning is a very rational decision.
    \item ACT-R fits best into the group of \texttt{humanly thinking}, because it deals with
    modeling human cognitive processes.
    \item Google Maps (i.e., its route planning tool) can be categorized as
    \texttt{rationally acting}, because it takes various factors into account, and
    based on that, the fastest route is calculated.
\end{enumerate}


\section*{Exercise 1.2}

\begin{enumerate}[(a)]
    \item First of all, we note that "scenic" is not a completely objective feature.
    Different people probably like different landscapes to varying degrees.
    Based on that, there might be ratings from previous customers about the route
    sections in a certain area. In such a scenario, the taxi might calculate the "most
    scenic" route with the highest "scenic value" by relying on these ratings.
    On the one hand, the taxi itself acts rationally, because it calculates the highest
    possible value. But on the other hand, the rating is human-related, so that would
    rather be \texttt{acting humanly}.
    \item The described behaviour of ChatGPT most likely is \texttt{acting rationally}.
    The reinforcement of human feedback has driven ChatGPT to more and more write in
    this way. This type of text (language) also is used by ChatGPT if it is not
    capable to give a substantive answer.
\end{enumerate}


\section*{Exercise 1.3}

\begin{enumerate}[(a)]
    \item The only possible yet very costly implementation that we could achieve is
    the following:
    For every step, check first all four directions. Then, among all non-blocked
    directions, choose randomly a direction and go there. If the goal was reached, stop.
    Else, check again and do next step.
    The problem is that the agent cannot remember where it was already. Another issue
    is that it cannot just follow a side (it does not "see" right and left).
    \item No, there is no guarantee that the agent will always reach the goal in
    finite many steps. Since the next step is always chosen randomly, it might end up never
    reaching the goal.
\end{enumerate}


\end{document}
